\section{Harmonic-flat approximation after a strict smooth boundary}
\label{sec:rpi:density-repair}

We pass from a general asymptotically flat end to a harmonically flat
one after constructing a strict smooth inner barrier. Uniform metric
comparison will control every enclosing area, while an annular flux
calculation controls the mass. Set \(d=n-2\), \(k=n-1\), \(\omega=|S^{n-1}|\), and
\(a_n=4k/d\).

\begin{lemma}[Smooth exterior harmonic-flat approximation]
\label{lem:rpi:smooth-density}
Let \((N,h)\) be a smooth connected asymptotically flat exterior,
complete including its compact smooth boundary \(\Gamma\), possibly
empty, with
one coordinate end and compact end complement. Suppose that
\[
 h-\delta=O_2(r^{-\tau}),\qquad \tau>d/2,\qquad
 R_h\ge0,\quad R_h\in L^1,\quad R_h=O(r^{-n-\beta})
\]
for some \(\beta>0\), and that its ADM mass exists.
There are smooth metrics \(\widehat h_R\), complete with the same
boundary, with nonnegative scalar curvature and a harmonically flat
end, such that
\[
 \begin{gathered}
 (1-\epsilon_R)h\le\widehat h_R\le(1+\epsilon_R)h,\qquad
 \epsilon_R\longrightarrow0,\\
 \widehat h_R\longrightarrow h\ \text{in }C^2
 \text{ on every fixed compact set},\\
 m_{\rm ADM}(\widehat h_R)\longrightarrow m_{\rm ADM}(h).
 \end{gathered}
\]
In particular, if \(H_\nu(\Gamma,h)<0\) for the normal toward the end,
the same strict sign holds for \(\widehat h_R\) when \(R\) is large.
\end{lemma}

\begin{proof}
We cut off the metric on a distant annulus and remove the resulting
scalar-curvature error by a Dirichlet conformal correction. The annular
flux then recovers the original ADM mass.

\smallskip\noindent\emph{Step 1: the cutoff and Dirichlet correction.}
Choose a smooth radial cutoff \(\chi_R=\chi(r/R)\), with
\(\chi=1\) on \((-\infty,1]\), \(\chi=0\) on \([2,\infty)\), and
\(0\le\chi\le1\). On the end put
\[
 h_R=\delta+\chi_R(h-\delta),
\]
and set \(h_R=h\) on the compact core. For large \(R\) this is a
smooth positive metric, equal to \(h\) on \(r\le R\) and Euclidean
on \(r\ge2R\). Extend \(\chi_R=1\) on the compact core and define
\[
 Q_R=R_{h_R}-\chi_R R_h.
\]
Thus \(Q_R\) is supported in the annulus
\(\mathcal A_R=\{R<r<2R\}\). The differentiated metric bounds and
the scalar-curvature formula give
\[
 |Q_R|\le C R^{-\tau-2},\qquad
 \|Q_R\|_{L^{n/2}(h_R)}\le C R^{-\tau},\qquad
 \|Q_R\|_{L^{2n/(n+2)}(h_R)}
       \le C R^{d/2-\tau}.
\]
The metrics \(h_R\) are uniformly comparable with \(h\), so the
Sobolev inequality
\[
 \|v\|_{L^{2^*}(h_R)}^2\le C\int_N|dv|_{h_R}^2\,dV_{h_R},
 \qquad 2^*=\frac{2n}{n-2},
\]
has a constant independent of large \(R\) on the homogeneous
completion of \(C_c^\infty(\operatorname{int}N)\). This is the
Euclidean-end Sobolev inequality combined with a compact-core
Poincar\'e inequality with trace on one fixed end annulus.
The completion imposes zero trace at \(\Gamma\).

Since \(\|Q_R\|_{n/2}\to0\), the bilinear form
\[
 \mathcal B_R(v,\varphi)
   =a_n\int\langle dv,d\varphi\rangle_{h_R}
      +\int Q_Rv\varphi
\]
is coercive, with
\(\mathcal B_R(v,v)\ge (a_n/2)\int|dv|^2\).
Solve
\[
 \mathcal B_R(v_R,\varphi)=-\int Q_R\varphi,\qquad z_R=1+v_R.
\]
Then
\[
 (-a_n\Delta_{h_R}+Q_R)z_R=0,\qquad
 z_R|_\Gamma=1,\qquad
 \|dv_R\|_2+\|v_R\|_{2^*}\le C R^{d/2-\tau}.
\]
All integrals in these estimates use \(h_R\).
Local elliptic regularity makes \(v_R\) smooth up to \(\Gamma\).
On the Euclidean end it is harmonic and belongs to \(L^{2^*}\);
the scaled mean-value estimate implies \(v_R\to0\) at infinity.

\smallskip\noindent\emph{Step 2: uniform and compact convergence.}
We first prove a uniform supremum estimate by annular rescaling. Cover \(\mathcal A_R\) by a fixed number of balls of
radius comparable to \(R\) in
\(\{R/2<r<4R\}\), and rescale each by \(R\). The rescaled metrics
have uniform ellipticity and bounded coefficients through order two;
the rescaled zeroth-order term \(R^2Q_R\) is \(O(R^{-\tau})\).
The scalar interior boundedness estimate for this equation gives
\[
 \sup_{\mathcal A_R}|v_R|
 \le C\left[
 R^{-d/2}\|v_R\|_{2^*}
       +R^2\|Q_R\|_\infty\right]
 \le C R^{-\tau}.
\]
The same bound holds on the closed source annulus by using these
slightly larger balls. Off that annulus \(v_R\) is harmonic, has
zero trace on \(\Gamma\), and tends to zero on the end. The maximum
principle on the inner and outer complementary regions therefore
gives \(\|v_R\|_\infty\le CR^{-\tau}\) on all of \(N\).
In particular \(z_R>0\) for large \(R\).
On each fixed compact region the equation for \(v_R\) is
\(\Delta_hv_R=0\) once \(R\) is large. Interior and fixed smooth
boundary estimates, with \(v_R=0\) on \(\Gamma\), give
\(v_R\to0\) in \(C^2\) there, and in every fixed higher norm.

Put \(\widehat h_R=z_R^{4/d}h_R\). The conformal scalar formula
now gives exactly
\[
 R_{\widehat h_R}
 =z_R^{-(n+2)/d}
       \bigl(-a_n\Delta_{h_R}z_R+R_{h_R}z_R\bigr)
 =z_R^{-4/d}\chi_R R_h\ge0.
\]
The global supremum estimate and
\(\sup_N|h_R-h|_h=O(R^{-\tau})\) give the stated metric comparison
and completeness. The end is \(z_R^{4/d}\delta\), with \(z_R\)
positive Euclidean harmonic and tending to one, hence is
harmonically flat to every order.

\smallskip\noindent\emph{Step 3: the ADM limit and the boundary flux.}
Write
\[
 z_R=1+A_Rr^{-d}+O_j(r^{-d-1})
\]
on that Euclidean end. Let \(\nu\) denote the \(h\)-unit normal at
\(\Gamma\) pointing into \(N\), and put
\(I_R=\int_\Gamma\partial_\nu z_R\,dA_h\).
The boundary \(C^2\) convergence gives \(I_R=o(1)\).
Since \(h_R=h\) near \(\Gamma\), integrating
\(a_n\Delta_{h_R}z_R=Q_Rz_R\) over a large truncation and then
letting its radius tend to infinity gives, with the outward
normal of the exterior domain equal to \(-\nu\) at \(\Gamma\),
\[
 -a_n d\omega A_R-a_n I_R
       =\int_N Q_Rz_R\,dV_{h_R}.
\]
The nonlinear factor is harmless:
\[
 \left|\int_NQ_Rv_R\,dV_{h_R}\right|
 \le\|Q_R\|_{2n/(n+2)}\|v_R\|_{2^*}
 \le C R^{d-2\tau}=o(1).
\]
For the remaining integral, write \(b=h-\delta\) and
\[
 \mathcal L(b)=\partial_i\partial_jb_{ij}
                         -\Delta_\delta\operatorname{tr}_\delta b.
\]
On \(\mathcal A_R\),
\[
 R_{h_R}\,dV_{h_R}
 =\left[\mathcal L(\chi_Rb)
                  +O(R^{-2\tau-2})\right]dx.
\]
Indeed all nonlinear scalar terms and the change of volume density
are bounded by \(C(|b_R||D^2b_R|+|Db_R|^2)\), where
\(b_R=\chi_Rb\). Their annular integral is \(O(R^{d-2\tau})\).
The outer linear flux is zero since \(b_R=0\) there; its inner flux
is the negative of the original ADM flux on \(S_R\). Consequently
\[
 \int_N Q_R\,dV_{h_R}
 =-\int_{S_R}(\partial_jh_{ij}-\partial_ih_{jj})n_\delta^i\,dA_\delta
     +O(R^{d-2\tau})
     -\int_{\mathcal A_R}\chi_RR_h\,dV_{h_R}.
\]
The last term tends to zero, by the scalar tail assumption
(or by integrability and uniform metric comparison). Thus
\[
 \int_NQ_R\,dV_{h_R}=-2k\omega m_{\rm ADM}(h)+o(1).
\]
As \(a_nd=4k\), the flux identity yields
\[
 A_R\longrightarrow \frac12m_{\rm ADM}(h),\qquad
 m_{\rm ADM}(\widehat h_R)=2A_R
                 \longrightarrow m_{\rm ADM}(h).
\]
Finally \(H_\nu(\Gamma,\widehat h_R)\to H_\nu(\Gamma,h)\)
uniformly by the compact \(C^2\) convergence. Strict negativity
therefore persists on the compact boundary.
\end{proof}

\begin{proof}[Proof of Theorem~\ref{thm:rpi}]
The case $3\le n\le7$ is the numerical Bray--Lee theorem described
in Section~\ref{sec:rpi:introduction}, with the classical positive
mass theorem when the frontier is empty. Assume $n\ge8$.
Proposition~\ref{prop:rpi:harmonic-flat-inequality} has proved the
nonempty-frontier case for harmonically flat ends. We now let $(E,h)$
have the general asymptotics of the theorem and full perimeter $A>0$.

\paragraph{The capacitary coefficient at infinity.}
Let \(w\) be its capacitary potential and \(c\) its capacity.
The local minimizing density bounds and the weak Harnack contraction
for the equilibrium potential, as in the proof of
Lemma~\ref{lem:rpi:flow-construction}, give its constant extension
the continuous boundary value one on the full frontier.
For any \(0<\alpha<\min\{\tau,1\}\), the positive
function \(r^{-d}(1-r^{-\alpha})\) is \(h\)-superharmonic for all
sufficiently large \(r\): its Euclidean Laplacian is
\(-\alpha(d+\alpha)r^{-n-\alpha}\), while the metric error is
\(O(r^{-n-\tau})\). Comparison with a fixed positive multiple on a
large coordinate sphere, followed by the maximum principle on the
end, gives \(w=O(r^{-d})\). Scaled interior estimates give
\(Dw=O(r^{-d-1})\), \(D^2w=O(r^{-d-2})\). Consequently
\[
 \Delta_\delta w
  =(\delta^{ij}-h^{ij})\partial_{ij}w
      +h^{ij}\Gamma_{ij}^\ell\partial_\ell w
  =O(r^{-n-\tau}).
\]
Extend a cutoff of \(w\), supported in its end chart, smoothly by zero
to \(\mathbb R^n\). Its Euclidean Laplacian has a finite
\(\alpha\)-moment. The Newton potential, split into inner, comparable
and outer annuli, is therefore a monopole plus
\(O_1(r^{-d-\alpha})\). Its difference from the extended function is
a decaying harmonic function on all of \(\mathbb R^n\), hence zero.
The cutoff has reduced this uniqueness statement to Euclidean space.
If the monopole coefficient is \(c_0\), flux integration on
\(\{w<s\}\), for regular values \(s\uparrow1\), gives
\[
 \int_{\{w<s\}}|dw|_h^2\,dV_h=s\,d\omega c_0.
\]
The flux at infinity is \(-d\omega c_0\); the other boundary is the
regular level \(w=s\), and the buried boundary is absent. Monotone
convergence and the capacity normalization identify \(c_0=c\).
Thus the finite-energy construction and this end argument give
\[
 0<w<1\ \text{in the open exterior},\qquad
 w=c r^{-d}+O_1(r^{-d-\alpha}),\qquad
 0<\alpha<\min\{\tau,1\}.
\]
In particular the conformal mass formula applies.

\paragraph{A strict smooth inner boundary.}
Fix a small \(\lambda>0\) and set
\[
 p=1+\lambda w,\qquad h_\lambda=p^{4/d}h.
\]
Then \(h_\lambda\ge h\), its asymptotically flat rate is
\(\min\{\tau,d\}>d/2\), and
\[
 R_{h_\lambda}=p^{-4/d}R_h\ge0,\qquad
 R_{h_\lambda}\,dV_{h_\lambda}=p^2R_h\,dV_h.
\]
Since $1\le p\le1+\lambda$, scalar integrability and the stated
scalar decay persist. The end expansion gives
\[
 m_{\rm ADM}(h_\lambda)=m_{\rm ADM}(h)+2\lambda c.
\]
To apply the static part of Corollary~\ref{cor:rpi:flow-detachment},
write this perturbation as
\[
 h^{\rm base}_\lambda=(1+\lambda)^{4/d}h,\qquad
 \widehat u_\lambda=\frac{1+\lambda w}{1+\lambda},\qquad
 h_\lambda=\widehat u_\lambda^{4/d}h^{\rm base}_\lambda.
\]
The normalized harmonic factor equals one on the frontier and lies
strictly between zero and one on the connected open exterior.
Constant scaling preserves local perimeter minimality of the original
filled set. Its smooth ambient metric supplies the minimizing-boundary
density and regular-patch estimates, and the harmonic factor has the
constant boundary data required by the separation argument. The
asymptotically flat end supplies the positive-mean-curvature coordinate
spheres used to confine the prescribed-curvature minimizer. Thus the
static detachment conclusion applies to the original conductor.
Proposition~\ref{prop:rpi:one-sided-smoothing} then gives a smooth
compact enclosing boundary $\Gamma$ with
$H_\nu(\Gamma,h_\lambda)<0$.
Retain its connected end exterior.
The metric \(h_\lambda\) is smooth on a neighborhood of its closure.
All parameters producing \(\Gamma\) are fixed before approximation.

\paragraph{Comparison of all enclosing areas and passage to the limit.}
Apply Lemma~\ref{lem:rpi:smooth-density} on this smooth exterior.
Let \(k_j\) be its approximating metrics, and choose
\(\epsilon_j\downarrow0\) so that
\[
 k_j\ge(1-\epsilon_j)h_\lambda,\qquad
 H_\nu(\Gamma,k_j)<0,\qquad
 m_{\rm ADM}(k_j)\longrightarrow m_{\rm ADM}(h_\lambda).
\]
Take a minimizing enclosure of the region removed by
\(\Gamma\) in \(k_j\). Lemma~\ref{lem:g-enclosure} gives compact existence, and
Lemma~\ref{lem:g-detachment} puts the entire minimizing frontier a
positive distance from the strict smooth obstacle. The resulting
frontier \(\Sigma_j\) is detached and locally perimeter minimizing,
and is outer minimizing with all components counted. Lemma~\ref{lem:g-enclosure} gives its connected end exterior. It has precisely the
boundary class of Theorem~\ref{thm:rpi}, a harmonically flat end and
scalar curvature zero on that distant end.

Every such enclosure contains the original conductor. Its full
perimeter in \(h\) is therefore at least \(A\), by the original
outer-minimizing property. This comparison applies directly to
finite-perimeter enclosures, or follows from strict smooth
approximation keeping a fixed neighborhood of the conductor.
Tangent-plane metric comparison gives
\[
 \mathcal H_{k_j}^{k}(\Sigma_j)
 \ge(1-\epsilon_j)^{k/2}
       \mathcal H_{h_\lambda}^{k}(\Sigma_j)
 \ge(1-\epsilon_j)^{k/2} A .
\]
Proposition~\ref{prop:rpi:harmonic-flat-inequality} consequently gives
\[
 m_{\rm ADM}(k_j)\ge
 \frac12\left(
       \frac{(1-\epsilon_j)^{k/2}A}{\omega}
       \right)^{d/k}.
\]
First let \(j\to\infty\) at fixed \(\lambda\), and then
\(\lambda\downarrow0\). The mass formula proves
\[
 m_{\rm ADM}(h)\ge
       \frac12(A/\omega)^{d/k}.
\]
The strict boundary and elliptic constants may depend on the fixed
\(\lambda\) and on \(\Gamma\); this is consistent with the stated
order of limits. The area estimate came from metric comparison and
conductor inclusion.

\paragraph{The empty frontier.}
A nonempty locally perimeter-minimizing frontier has
positive perimeter by the density lower bound. If \(A=0\), the
frontier is therefore empty. On that complete boundaryless manifold
apply the same cutoff-density proof with \(\Gamma=\varnothing\);
there is then no compact-boundary flux. Applying
Lemma~\ref{lem:rpi:harmonic-pmt} to the harmonic-flat approximants
and passing their masses to the limit gives \(m_{\rm ADM}(h)\ge0\).
\end{proof}
